\documentclass[UTF8,fontset=none,11pt,a4paper]{ctexart}
\usepackage{ipara-handbook}
\usepackage{amsmath,booktabs,tabularx,hyperref}
\usetikzlibrary{arrows.meta,positioning,fit,calc}

\handbooksetup{OS-I01 一次 Blocking read()}{408 · Operating System Integration}{Task · File Reference · Cache · I/O · Wakeup}

\begin{document}
\handbooktitle{一次 Blocking \texttt{read()} 的一生}{把进程、文件、页缓存、I/O 与调度组合成一条可验证状态轨迹}{408 · OS-I01}

\begin{corebox}{Canonical Problem}
一个正在运行的进程调用 \code{read(fd, buf, n)}。若所需文件内容不能立即从内存状态满足，系统怎样从用户态进入内核，沿文件引用找到数据，把存储请求交给设备，让当前 task 安全等待，并在完成后重新获得执行机会，最终返回“实际读取字节数 / EOF / error”？
\[
\boxed{
\text{Running Task}
\to \text{Syscall}
\to \text{fd/OFD/File Range}
\to \text{Cached/Resident?}
\to \text{I/O if needed}
\to \text{Block/Wake}
\to \text{Resume/Return}
}
\]
箭头表示本次组合过程中的依赖与状态推进。Cache hit、EOF、error 等分支可以提前结束慢路径，但不能跳过相应合法性检查。
\end{corebox}

\begin{methodbox}{Integration Ownership}
本册只 Owns \kw{模块识别、组合次序、跨模块状态轨迹、分支与独立验证}。
\begin{itemize}
  \item OS-01/02 Own task、syscall 后的控制权、Blocked/Ready、scheduler/dispatcher；
  \item OS-06/07 Own fd/OFD/file object、current offset、file range 与 byte-to-block；
  \item OS-B03 Own process descriptor binding 与 open instance 的 handoff；
  \item OS-B04 Own file range、cached/resident page、VA mapping 与 I/O 之间的 handoff；
  \item OS-05 Own request、driver、device progress、completion；
  \item OS-B01 Own wait relation、block、wake、重新竞争 CPU 的接口；
  \item 若题目继续追问 interrupt/DMA 的硬件交接，再调用 X-B03。
\end{itemize}
本 Integration 不重新定义这些机制。
\end{methodbox}

\section{初始状态、目标与四条并行轨迹}

题目一旦出现 \code{read()}，不要只画一条“进程阻塞”的线。至少同时追踪四条状态：

\begin{handbooklongtable}{L{2.6cm} Y Y Y}
\toprule
轨迹 & 初始问题 & 关键状态 & 成功时输出\\
\midrule
Task / Control & 谁正在执行？CPU 控制权在哪？ & user/kernel、Running/Blocked/Ready、wait relation & task 最终在某次调度后继续，并按 syscall ABI 返回\\
File Reference & fd 指向谁？当前文件位置是什么？ & fd binding、OFD/open instance、file object、offset & 定位本次实际读取的 file range\\
Data / Residency & 所需内容现在在哪里？ & user buffer、cached file page、resident/dirty/I/O-in-flight & 目标 bytes 已可从内存状态交付给本次 read\\
Device / Completion & 若内存不足以满足，谁负责推进存储请求？ & request queue、driver/device state、completion evidence & 请求完成状态被软件确认，并可唤醒 waiter\\
\bottomrule
\end{handbooklongtable}

成功条件不是“磁盘读完了”，也不是“进程被唤醒了”，而是：\kw{本次 read 的接口语义得到一个合法返回结果，进程在以后真正获得 CPU 后把该结果返回用户态。}

\section{阶段一：从用户调用进入内核，但暂时还是同一个 task}

进程在用户态执行 \code{read(fd, buf, n)}，通过系统调用入口把控制权交给内核。此时首先发生的是 privilege/control transfer，不应自动画成进程切换：
\[
\boxed{\text{Mode/Privilege Entry} \not\Rightarrow \text{Context Switch}}
\]

内核先验证参数和当前调用语义所需的权限/对象条件。典型检查包括：fd 是否有效、打开实例是否允许读取、用户 buffer 是否可合法访问、请求长度是否合法。具体 ABI 与 user-pointer 校验机制由 OS-00/OS-01/02 的接口边界解释。

\section{阶段二：沿 fd 找到“这次读的是文件的哪一段”}

通过 OS-B03 的稳定关系：
\[
\boxed{\text{Task}\xrightarrow{fd}\text{Open Instance/OFD}\to\text{File Object}}
\]
内核得到本次操作依赖的打开实例状态，并确定逻辑读取起点。对普通顺序读，current offset 与“实际成功消费的 bytes”有关；如果题目使用 \code{pread()} 一类显式位置接口，则不能再假设共享 OFD offset 被推进。

于是文件系统把问题压成：
\[
\boxed{\text{File Object} + [offset,\ offset+n)\text{ 的实际可读部分}}
\]
EOF、文件长度、设备/特殊文件语义都可能让本次调用在产生存储 I/O 之前就得到结果。

\section{阶段三：先问内存状态能不能立即满足，而不是先画磁盘}

OS-B04 要求先把\kw{文件内容身份}与\kw{驻留状态}分开：

\begin{enumerate}
  \item \textbf{Cached/Resident Hit}：所需内容已有可用内存副本，本次 read 可以直接进入数据交付路径；通常不需要因“文件在磁盘上”而把 task 阻塞等待设备。
  \item \textbf{Content Not Resident / Need More Data}：现有内存状态不足以满足本次读取，文件系统确定需要哪些后备存储数据，并向 OS-05 形成 I/O request。
\end{enumerate}

\begin{warnbox}{第一个关键 Anti-Bridge}
\[
\boxed{\text{read(file)} \not\Rightarrow \text{disk I/O}}
\]
文件是持久对象，不代表每次访问都必须碰设备。Page Cache / buffer hit 可以让调用只走内存快路径；反过来，cache miss 也不自动等于“必须置换页面”，还要看是否有可用 frame 和具体分配状态。
\end{warnbox}

\section{阶段四：真正需要 I/O 时，把“请求推进”和“task 等待”拆成两条线}

当数据必须由设备推进时：
\begin{enumerate}
  \item OS-06/07 把 file range 翻译成下游可请求的数据位置/块关系；
  \item OS-05 建立 request，交给 driver/controller/device；
  \item 如果当前接口语义是 blocking 且结果尚不能返回，OS-B01 建立“当前 task 正在等待这个完成条件”的 wait relation，并让 task 失去 CPU 调度资格；
  \item scheduler 从其他 Ready/Runnable 实体中选择下一位，必要时发生 context switch。
\end{enumerate}

这里必须守住：
\[
\boxed{\text{Submit I/O} \neq \text{Block}}
\]
异步接口可以提交后先返回；即使 blocking read，也可能因已有数据、partial result 或 error 而无需等待。\kw{是否 block 是调用者/接口状态，谁搬运数据是 I/O mechanism，两者是正交维度。}

\section{阶段五：设备完成先改变软件状态，再改变 task 的调度资格}

设备请求在 CPU 之外继续推进。完成以后，系统必须先得到足够的 completion evidence，并把结果关联回正确 request：
\[
\text{Device Progress}
\to \text{Completion Detected}
\to \text{Request/Page State Updated}
\to \text{Wake Eligible Waiter}
\]

完成发现可能依赖 interrupt、polling 或混合方案；DMA 可能参与数据搬运，也可能不是题目所在抽象层。Integration 不把某一种硬件实现当成定义。

最关键的状态边界：
\[
\boxed{\text{I/O Completion} \to \text{Ready/Runnable} \not\to \text{Running}}
\]
被唤醒的 task 只是重新进入 CPU 竞争集合。是否立即抢占当前运行者由调度策略、安全点与优先级等条件决定。

\section{阶段六：重新运行以后，系统调用从“旧内核路径”继续}

当调度器以后真正选中原 task：
\begin{enumerate}
  \item 恢复该 task 的可恢复执行现场；
  \item 重新检查本次 request / cached content 是否满足继续条件；
  \item 把实际可交付数据推进到调用者 buffer（具体是否发生用户--内核复制取决于路径与接口）；
  \item 更新本次操作拥有的文件位置/返回状态；
  \item 通过 syscall return 恢复到用户态，返回实际字节数、0（例如 EOF）或 error。
\end{enumerate}

注意 \code{read(fd,buf,n)} 的成功返回值不承诺一定等于 $n$。题目若涉及 pipe、terminal、socket、signal interruption 或 partial I/O，应按对应对象接口重新判断；本册的 canonical mother example 是普通 file-backed blocking read。

\section{完整母轨迹：Cache Miss 的普通文件读取}

\begin{center}
\small
\begin{tikzpicture}[node distance=7mm and 7mm,>=Stealth]
\tikzset{n/.style={draw=iparaDeep,fill=iparaSoft,rounded corners=1.5mm,align=center,minimum width=2.55cm,minimum height=0.78cm,font=\scriptsize}}
\node[n] (u) {Task Running\\user \code{read()}};
\node[n,right=of u] (k) {Kernel Entry\\validate fd/buf};
\node[n,right=of k] (f) {fd $\to$ OFD\\$\to$ file range};
\node[n,below=of f] (miss) {Content not resident\\form I/O request};
\node[n,left=of miss] (block) {Register wait\\Blocked};
\node[n,left=of block] (other) {Other task runs\\device progresses};
\node[n,below=of other] (done) {Completion detected\\data state updated};
\node[n,right=of done] (ready) {Wake\\Ready/Runnable};
\node[n,right=of ready] (resume) {Scheduler/Dispatch\\resume syscall};
\node[n,above=of resume] (ret) {deliver bytes\\return user};
\draw[->,thick] (u)--(k); \draw[->,thick] (k)--(f); \draw[->,thick] (f)--(miss);
\draw[->,thick] (miss)--(block); \draw[->,thick] (block)--(other); \draw[->,thick] (other)--(done);
\draw[->,thick] (done)--(ready); \draw[->,thick] (ready)--(resume); \draw[->,thick] (resume)--(ret);
\end{tikzpicture}
\end{center}

这张图里只有两条地方真正代表 CPU execution context change：task block 后换给别人，以及以后 scheduler 重新选择原 task。Syscall entry/return 是 privilege transition；device completion 是事件；wake 是资格变化。把四者混成“上下文切换”会直接破坏计算题计数。

\section{高价值分支：同一句 read，轨迹为什么不同}

\begin{handbooklongtable}{L{3.0cm} Y Y}
\toprule
条件变化 & 轨迹变化 & 不应错误推出\\
\midrule
Page Cache hit & 跳过设备 request 与 I/O wait，直接交付可用数据 & “read 一定阻塞”\\
到达 EOF & 可直接返回 0 或已读取的 partial bytes（按接口状态） & “必须读满 $n$ 才返回”\\
无空闲 frame & VM 可能先需要 reclaim/replacement，再形成/完成 page-in & “cache miss 就等于 replacement”\\
I/O 完成时 CPU 正跑别的 task & waiter 先 Ready，是否抢占再由 scheduler 决定 & “interrupt 一来原 task 立即继续”\\
同一 OFD 被 fork/dup 共享 & 成功读取消费的 current offset 可能影响另一引用者 & “不同 fd 数字就有独立 offset”\\
使用 \code{pread()} & 显式位置读取，不按普通 current offset 推进 & “任何 read-like API 都共享 OFD offset”\\
使用 \code{mmap()} & 文件内容通过 VA mapping + fault/load 访问，不再以每次 read syscall 为主线 & “mmap 只是更快的 read”\\
Direct/async I/O & cache/copy/blocking 关系可能改变 & “DMA/async/nonblocking 是同义词”\\
\bottomrule
\end{handbooklongtable}

\section{失败定位：找第一个不能满足的前置条件}

\begin{handbooktable}{L{3.0cm} Y Y}
\toprule
现象 & 首个候选 Owner & 第一验证动作\\
\midrule
立即返回 bad fd / permission error & OS-06/07 + OS-B03 & fd binding、打开模式、file object 是否合法\\
地址/用户 buffer 错误 & OS-01/02 + OS-04 Use & user pointer 的映射/权限是否允许当前 copy\\
文件存在但读到 EOF & OS-06/07 & current offset 与 file size/range\\
长时间阻塞 & OS-05 + OS-B01 & request 是否提交、等待哪个 completion、是否已 wake\\
设备完成但进程未运行 & OS-01/02 & task 是否 Ready、scheduler 是否已选择它\\
读到内容但 offset 推演错 & OS-B03 & 是否共享同一 OFD、实际返回多少 bytes\\
\bottomrule
\end{handbooktable}

\section{独立验证与停止条件}

完成一道 blocking read 综合题后，固定做六次检查：
\begin{enumerate}
  \item \kw{Object}：fd、OFD、file object、file range、cached page、I/O request、task 是否被当成不同对象？
  \item \kw{Control}：syscall entry、block/context switch、interrupt/completion、wake、dispatch 是否分开？
  \item \kw{Queue/Relation}：task 真正因什么条件进入哪个 wait relation，什么事件解除它？
  \item \kw{Data}：数据此刻到底在 persistent storage、cache/resident memory 还是 user buffer？
  \item \kw{Accounting}：offset 只按实际消费的 bytes 与具体接口契约变化，是否误把请求长度 $n$ 当实际完成量？
  \item \kw{Stop}：如果题目没问 DMA/MMU/文件系统持久化细节，就在对应 Bridge 接口停止，不复制 Topic 内部机制。
\end{enumerate}

\begin{corebox}{Compression}
遇到“阻塞读”不要背一条固定故事，按四问复原：
\[
\boxed{
\text{读哪一个 file range？}
\to \text{内存能否立即满足？}
\to \text{若不能，谁推进 request、task 等什么？}
\to \text{completion 后为什么只是 Ready，何时才 return？}
}
\]
能回答这四问，OS-01/02、OS-05、OS-06/07 与 B01/B03/B04 就已经正确组合起来。
\end{corebox}

\end{document}
